% \documentclass[10pt,a4paper]{article}
%
% \usepackage[utf8]{inputenc}
% \usepackage[T1]{fontenc}
% \usepackage[margin=1.7cm,top=1.4cm,bottom=1.5cm]{geometry}
% \usepackage{booktabs}
% \usepackage{array}
% \usepackage{amsmath}
% \usepackage{amssymb}
% \usepackage{xcolor}
% \usepackage{colortbl}
% \usepackage{enumitem}
% \usepackage{titlesec}
% \usepackage[hidelinks]{hyperref}
% \usepackage{caption}
%
% \definecolor{headblue}{RGB}{20,60,110}
% \definecolor{rowgray}{RGB}{242,244,248}
%
% \titleformat{\section}{\large\bfseries\color{headblue}}{\thesection}{0.5em}{}
% \titlespacing*{\section}{0pt}{8pt}{4pt}
% \titleformat{\subsection}{\normalsize\bfseries\color{headblue}}{\thesubsection}{0.5em}{}
% \titlespacing*{\subsection}{0pt}{6pt}{2pt}
%
% \setlength{\parskip}{2pt}
% \setlength{\parindent}{0pt}
% \captionsetup{font=footnotesize,labelfont=bf,skip=3pt}
%
% \newcommand{\lam}[3]{$\lambda_{\mathrm{SR}}{=}#1,\ \lambda_{\mathrm{IR}}{=}#2,\ \lambda_{\mathrm{SU}}{=}#3$}
%
% \title{\vspace{-1.6cm}\textbf{Task-Arithmetic Model Merging for the \emph{clearsky} Project}\\[2pt]
% \normalsize Condensed report: evolution of the merging search and its selection metric}
% \author{}
% \date{}
%
% \begin{document}
% \maketitle
% \vspace{-1.4cm}
%
% \section*{Setup}
% \emph{clearsky} restores degraded astronomical images with a single conditional
% diffusion model (U-Net with self-attention, GroupNorm/SiLU, sinusoidal time
% embeddings), trained as a v-prediction DDPM with a cosine schedule. Three
% single-task experts --- \textbf{star removal}, \textbf{image restoration}, and
% \textbf{super-resolution} --- were fine-tuned from one shared base $\theta_0$.
% Task arithmetic merges them without retraining: each expert gives a task vector
% $\tau_t=\theta_t-\theta_0$, combined as
% \[
% 	\theta_{\mathrm{merged}} = \theta_0 + \sum_{t}\lambda_t\,\tau_t ,
% \]
% with per-task coefficients $\lambda_t$ chosen by a validation search. Task-vector
% norms were stable across all runs
% ($\lVert\tau_{\mathrm{SR}}\rVert{\approx}56.4$,
% $\lVert\tau_{\mathrm{IR}}\rVert{\approx}50.3$,
% $\lVert\tau_{\mathrm{SU}}\rVert{\approx}80.9$): super-resolution moves far
% further from the base than the other two, an early hint that it would dominate
% or interfere in a naive combination. A folder/task naming mismatch (which
% abbreviation meant which task) also had to be tracked down and corrected along
% the way.
%
% \section*{Evolution of the search and, more importantly, the metric}
% The real story here is not the search algorithm but the \emph{objective it
% 	optimised}.
%
% \textbf{Raw, pooled MSE.} The earliest attempts scored candidate merges by
% mean-squared error against ground truth, first at a single hand-picked
% $\lambda$ point and then over an actual grid, split into a two-stage
% search (an easy ``safe zone'' of single-degradation folders, followed by
% harder combined-degradation folders). Both consistently picked
% $\lambda_{\mathrm{SU}}=0$: because the super-resolution folder carries by far
% the largest absolute MSE, any pooled objective is minimised by
% \emph{suppressing} that task rather than balancing it --- a metric artifact,
% not a genuine result, and unaffected by how much richer the grid or the
% two-stage design became.
%
% \textbf{Normalised MSE.} The fix was to measure each task's error relative to
% its own single-task expert,
% \[
% 	\mathrm{score}=\frac{1}{|T|}\sum_{t}\frac{\mathrm{MSE}_{\mathrm{merge}}[t]}{\mathrm{MSE}_{\mathrm{expert}}[t]},
% \]
% computed with one loader per task and a fixed noise seed so every comparison is
% fair, alongside PSNR/SSIM for a fuller picture. Folded back into the two-stage
% grid (after fixing the folder$\to$task label bug), this normalised objective
% was the first to select a genuine three-task compromise,
% \textbf{\lam{0.9}{0.1}{0.1}}, instead of switching a task off entirely.
% Combined-degradation folders have no single expert to normalise against, so
% they were instead scored relative to $\theta_0$ --- i.e.\ how much merging
% improves on doing nothing.
%
% \textbf{Perceptual Z-score.} The final refinement replaced MSE with a
% perceptual metric (LPIPS) on any folder involving super-resolution, and put
% every folder on a common scale via a per-folder Z-score
% ($z=(x-\mu)/\sigma$ across the candidate grid, averaged) before ranking. This
% was prompted by a mismatch between low MSE and visibly poor super-resolution
% output quality.
%
% \begin{table}[h]
% 	\centering
% 	\footnotesize
% 	\renewcommand{\arraystretch}{1.3}
% 	\caption{Evolution of the search objective. $\lambda$ reported as (SR/IR/SU).}
% 	\begin{tabular}{@{}p{3.1cm}p{1.6cm}p{2.0cm}p{6.6cm}@{}}
% 		\toprule
% 		\textbf{Selection metric}            & \textbf{$\lambda$ grid} & \textbf{Winner}      & \textbf{Key result}                             \\
% 		\midrule
% 		\rowcolor{rowgray}
% 		Raw MSE, single point                & 1 pt                    & 0.9/0.1/0.0          & SU switched off entirely                        \\
% 		Raw MSE, two-stage grid              & 27 combos               & 1.0/0.1/0.0          & SU still off; loss dominated by SU folder       \\
% 		\rowcolor{rowgray}
% 		Normalised MSE (proposed)            & $\le$4 combos           & truncated run        & establishes fair, scale-free per-task metric    \\
% 		Normalised MSE, two-stage, bug fixed & 36 combos               & \textbf{0.9/0.1/0.1} & first genuine three-task compromise             \\
% 		\rowcolor{rowgray}
% 		Perceptual Z-score (LPIPS + MSE)     & 8 combos                & 0.9/0.1/0.1          & same $\lambda$, but exposes hidden interference \\
% 		\bottomrule
% 	\end{tabular}
% \end{table}
%
% \section*{What the diagnostics actually showed}
% Even though the last two objectives agree on \lam{0.9}{0.1}{0.1}, per-task
% diagnostics on the super-resolution folder show this merge is not the clean win
% it appears to be:
%
% \begin{table}[h]
% 	\centering
% 	\footnotesize
% 	\renewcommand{\arraystretch}{1.25}
% 	\caption{Super-resolution folder: deterministic val.\ loss, MSE, LPIPS.}
% 	\begin{tabular}{@{}lccc@{}}
% 		\toprule
% 		\textbf{Model}                      & \textbf{Det.\ val loss} & \textbf{MSE}     & \textbf{LPIPS}  \\
% 		\midrule
% 		Pure pre-trained base $\theta_0$    & 0.0835                  & 0.00827          & 0.2303          \\
% 		Single-task super-resolution expert & 0.0880                  & 0.00893          & 0.2203          \\
% 		Merged model (0.9/0.1/0.1)          & 0.1557                  & \textbf{0.00707} & \textbf{0.4080} \\
% 		\bottomrule
% 	\end{tabular}
% \end{table}
%
% The merge has the \emph{lowest} MSE of the three yet by far the \emph{worst}
% LPIPS: MSE was rewarding a smoother, perceptually worse prediction. The LPIPS
% gap persists even with many more sampling steps, ruling out a sampling
% artifact. The task vectors are nearly orthogonal (cosine similarity
% $\approx 0.08$--$0.12$), which normally favours task arithmetic, but
% $\tau_{\mathrm{SU}}$'s much larger norm means even a small mixing coefficient
% from the other two tasks perturbs its direction enough to visibly degrade
% perceptual quality --- genuine task-vector interference, not noise.
%
% \section*{Conclusion}
% The whole arc reduces to one lesson: \textbf{the choice of evaluation metric
% 	mattered far more than the choice of search algorithm.} Raw pooled MSE quietly
% answered a different question --- ``minimise total pixel error'' --- than the
% one actually being asked --- ``be good at all three tasks at once'' --- which
% is why it kept switching super-resolution off. Normalising against per-task
% baselines fixed that pooling artifact and recovered a real compromise, but only
% a perceptual metric exposed the deeper issue: interference between task
% vectors that pixel-wise error cannot see. Any merge of these experts should be
% validated perceptually before being trusted; promising next steps are
% interference-aware merging (e.g.\ TIES/DARE), renormalising the
% super-resolution task vector, or building it from an earlier, smaller-norm
% checkpoint.
%
% \end{document}

\documentclass[11pt,a4paper]{article}

% ----------------------------------------------------------------------
% Packages
% ----------------------------------------------------------------------
\usepackage[utf8]{inputenc}
\usepackage[T1]{fontenc}
\usepackage[margin=2.4cm]{geometry}
\usepackage{booktabs}
\usepackage{array}
\usepackage{longtable}
\usepackage{amsmath}
\usepackage{amssymb}
\usepackage{xcolor}
\usepackage{colortbl}
\usepackage{graphicx}
\usepackage{enumitem}
\usepackage{titlesec}
\usepackage[hidelinks]{hyperref}
\usepackage{caption}

\definecolor{headblue}{RGB}{20,60,110}
\definecolor{rowgray}{RGB}{242,244,248}
\definecolor{codebg}{RGB}{246,246,246}

\titleformat{\section}{\Large\bfseries\color{headblue}}{\thesection}{0.6em}{}
\titleformat{\subsection}{\large\bfseries\color{headblue}}{\thesubsection}{0.6em}{}

\newcommand{\code}[1]{\texttt{\small #1}}
\newcommand{\lam}[3]{$\lambda_{\mathrm{SR}}{=}#1,\ \lambda_{\mathrm{IR}}{=}#2,\ \lambda_{\mathrm{SU}}{=}#3$}

\captionsetup{font=small,labelfont=bf}

% ----------------------------------------------------------------------
\title{\vspace{-1.2cm}\textbf{Task-Arithmetic Model Merging for the \emph{clearsky} Project}\\[4pt]
\large A Historical Log of the Merging Experiments, Their Rationale and Findings}
\author{Experimental report}
\date{\today}

\begin{document}
\maketitle
\vspace{-0.6cm}

% ======================================================================
\section{Context and Setup}

The \emph{clearsky} project restores degraded astronomical images with a
conditional diffusion model. A single \code{ImprovedUNet} (a U-Net with
self-attention, sinusoidal time embeddings, GroupNorm/SiLU blocks) is trained as a
\textbf{v-prediction DDPM} with a cosine noise schedule. Three separate
single-task \emph{experts} were fine-tuned from one shared pre-trained
base $\theta_0$:

\begin{itemize}[nosep,leftmargin=1.4em]
   \item \textbf{star\_removal} --- removal of stars from the field;
   \item \textbf{image\_restoration} --- repair of dead / hot pixels;
   \item \textbf{super\_resolution} --- upsampling / detail recovery.
\end{itemize}

The goal of every experiment below is the same: obtain \emph{one} network that
handles all three degradations at once, without retraining, using
\textbf{task arithmetic}. For each expert one computes a \emph{task vector}
$\tau_t=\theta_t-\theta_0$, and the merged weights are

\[
\theta_{\mathrm{merged}} \;=\; \theta_0 \;+\; \sum_{t}\lambda_t\,\tau_t ,
\]

where the per-task coefficients $\lambda_t$ are chosen by a search over a
validation set. The five notebooks are a chronological record of how the
\emph{search strategy} and, more importantly, the \emph{selection metric}
were refined as flaws in earlier choices became visible. The network
architecture, the v-prediction objective and the merging formula itself never
change; everything that follows is about \emph{how the $\lambda$'s are chosen
and judged}.

\paragraph{A naming caveat carried through the log.}
The dataset folders and the weight names do \emph{not} share the same
abbreviations, which caused a real bug that experiment~4 had to fix. The
convention finally adopted is: folder \code{SR}\,$\rightarrow$\,star\_removal,
folder \code{SU}\,$\rightarrow$\,super\_resolution, folder
\code{IR}\,$\rightarrow$\,image\_restoration. Combined-degradation folders
(\code{SR\_IR}, \code{SR\_SU}, \code{IR\_SU}, \code{SR\_IR\_SU}) hold images
corrupted by several degradations simultaneously.

Throughout, the task vector norms were stable across runs
($\lVert\tau_{\mathrm{SR}}\rVert\!\approx\!56.4$,
$\lVert\tau_{\mathrm{IR}}\rVert\!\approx\!50.3$,
$\lVert\tau_{\mathrm{SU}}\rVert\!\approx\!80.9$), i.e.\ the super-resolution
expert moved much further from the base than the other two --- an early hint
that it would dominate any naive combination.

% ======================================================================
\section{Summary of the Five Experiments}

Table~\ref{tab:summary} condenses the whole progression. The central story is
the migration of the \emph{selection metric}: from raw MSE (Exp.~1--2), to
MSE normalised against per-task baselines (Exp.~3--4), to a perceptual,
standardised Z-score (Exp.~5). Each change was forced by a specific failure of
the previous metric.

\renewcommand{\arraystretch}{1.35}
\begin{table}[h]
\centering
\footnotesize
\caption{Historical evolution of the merging experiments. Winner reported as
\lam{}{}{} coefficients. ``score'' meaning differs per row (see the metric column).}
\label{tab:summary}
\begin{tabular}{@{}p{0.9cm} p{2.5cm} p{3.4cm} p{2.4cm} p{2.0cm} p{2.2cm}@{}}
\toprule
\textbf{Exp.} & \textbf{Search strategy} & \textbf{Selection metric} &
\textbf{$\lambda$ grid} & \textbf{Winner (SR/IR/SU)} & \textbf{Headline result} \\
\midrule
\rowcolor{rowgray}
1 & Single point (confirmation run) &
Raw MSE, pooled per-task val splits &
1 point &
$0.9/0.1/0.0$ &
val\_loss $=0.01237$; SU dropped \\

2 & Two-stage grid (safe zone $\to$ fine balance) &
Raw MSE per folder, averaged &
27 combos &
$1.0/0.1/0.0$ &
val\_loss $=0.00398$; SU still $0$ \\

\rowcolor{rowgray}
3 & Single-stage grid &
\textbf{Normalised} MSE/expert; adds PSNR, SSIM &
$\le$4 combos &
(run truncated) &
Introduces fair, scale-free metric \\

4 & Two-stage grid; folder$\to$task bug fixed &
Norm.\ vs \emph{expert} (stage 1) \& vs \emph{$\theta_0$} (stage 2) &
36 combos &
$0.9/0.1/0.1$ &
score $=0.1091$; \textbf{SU now non-zero} \\

\rowcolor{rowgray}
5 & Two-stage grid + diagnostics &
\textbf{Z-score}: LPIPS on SU, MSE on SR/IR &
8 combos &
$0.9/0.1/0.1$ &
Reveals perceptual interference on SU \\
\bottomrule
\end{tabular}
\end{table}

% ======================================================================
\section{Experiment 1 --- Baseline Merge and Raw-MSE Evaluation}

\subsection*{What it does}
The first notebook establishes the full pipeline: load $\theta_0$ and the three
expert checkpoints (from the \code{ema\_model} key), check shape/key
compatibility, build the task vectors, and evaluate a merged model on isolated
per-task validation splits selected through a \code{split\_manifest} JSON.
Evaluation samples are produced with a fast \textbf{DDIM} sampler
(25 steps) and scored by \textbf{mean squared error} against ground truth;
the per-task MSEs are pooled into a single \code{val\_loss}. A final,
full-quality \textbf{DDPM} pass re-checks the chosen model.

\subsection*{Rationale}
This is the minimum viable merging loop --- it answers ``does the arithmetic
even produce a usable network?'' As shipped, the $\lambda$ grid holds a
\emph{single} point, \lam{0.9}{0.1}{0.0}: the notebook is a confirmation run of a
coefficient set found in an earlier exploratory sweep (a code comment records a
prior best of \lam{0.5}{0.6}{0.0} at val\_loss $0.00553$).

\subsection*{Findings}
The merged model reaches \code{val\_loss} $=0.012368$, with a very uneven
per-task breakdown: star\_removal MSE $\approx 7\times10^{-5}$ but
image\_restoration $\approx 0.0268$ and super\_resolution $\approx 0.0222$.
\textbf{Super-resolution is switched off entirely} ($\lambda_{\mathrm{SU}}=0$).
This is the first symptom of the problem that dominates the rest of the log:
because SU carries by far the largest absolute MSE, any pooled-MSE objective is
minimised by \emph{suppressing} SU rather than by balancing the three tasks.

% ======================================================================
\section{Experiment 2 --- Two-Stage Grid Search}

\subsection*{What changed}
\begin{itemize}[nosep,leftmargin=1.4em]
   \item A genuine \textbf{grid search} replaces the single point:
         $\lambda_{\mathrm{SR}}\in\{0.8,0.9,1.0\}$,
         $\lambda_{\mathrm{IR}}\in\{0.1,0.3,0.5\}$,
         $\lambda_{\mathrm{SU}}\in\{0.0,0.1,0.3\}$ --- 27 combinations.
   \item The search is split into \textbf{two stages}. \emph{Stage 1}
         (``safe zone'') scores every combination on the single-degradation
         folders \code{IR/SR/SU} and keeps the best 10. \emph{Stage 2}
         re-scores only those 10 on the combined-degradation folders
         (\code{SR\_IR}, \code{SR\_SU}, \code{IR\_SU}, \code{SR\_IR\_SU}) to pick
         the final balance.
   \item The dataset class is rewritten to pool several roots and take a
         \textbf{deterministic random subsample} (5\% per folder), because the
         full sets are too large; the deterministic seed keeps every DDP process
         aligned on the same slice.
\end{itemize}

\subsection*{Rationale}
Exp.~1 only ever tested hand-picked points. The two-stage design encodes a
reasonable intuition: first find coefficients that each task tolerates in
isolation (a ``safe zone''), then break ties using the harder, realistic
combined-degradation images --- while keeping the combinatorics affordable by
subsampling.

\subsection*{Findings}
Stage~1's best is \lam{0.9}{0.1}{0.0} (val\_loss $0.002975$); stage~2's best is
\lam{1.0}{0.1}{0.0} (val\_loss $0.003978$). The machinery is richer, but the
\textbf{metric is still raw averaged MSE}, so the outcome is the same pathology
as Exp.~1: \textbf{$\lambda_{\mathrm{SU}}=0$ wins again}. The averaged loss is
dominated by the \code{SU} folder (its MSE $\sim 0.006$ dwarfs \code{IR}/\code{SR}
at $\sim 0.001$), so the optimiser is rewarded for ignoring super-resolution.
This makes explicit that the \emph{search} was never the bottleneck --- the
\emph{objective} was.

% ======================================================================
\section{Experiment 3 --- A Fair, Normalised Metric}

\subsection*{What changed}
\begin{itemize}[nosep,leftmargin=1.4em]
   \item \textbf{Normalised selection metric.} Instead of raw MSE, the score is
         \[
           \mathrm{score}=\frac{1}{|T|}\sum_{t}\frac{\mathrm{MSE}_{\mathrm{merge}}[t]}
           {\mathrm{MSE}_{\mathrm{expert}}[t]},
         \]
         i.e.\ each task's error is measured \emph{relative to its own expert}.
         This is invariant to the very different MSE scales/floors of the three
         tasks and no longer collapses onto whichever task has the largest
         absolute error.
   \item \textbf{Expert baselines} are computed once, up front (each expert
         evaluated only on its own task).
   \item \textbf{One loader per task} (no pooling), and a \textbf{fixed per-task
         seed}: the same initial noise is used for the expert and for every
         $\lambda$ combination, so comparisons are fair.
   \item \textbf{Perceptual metrics} PSNR and a dependency-free \textbf{SSIM}
         are added to the report; results are dumped to JSON.
\end{itemize}

\subsection*{Rationale}
This is the conceptual turning point. Exp.~1--2 proved that averaging raw MSE
systematically discards the highest-error task. Normalising each task by its own
expert puts the three tasks on a common footing --- a merge that reaches, say,
$1.2\times$ its expert's error on every task is now clearly preferable to one
that is perfect on two tasks and absent on the third. PSNR/SSIM were added
because MSE alone says little about perceived image quality.

\subsection*{Findings}
The logged run is truncated (only the \code{image\_restoration} expert baseline
prints before the notebook moves on), so this experiment contributes the
\emph{method} rather than a final winner. Its ideas --- normalisation, one-off
expert baselines, fixed seeds, perceptual reporting --- are inherited by every
later experiment.

% ======================================================================
\section{Experiment 4 --- Two-Stage Search with Normalised Metrics}

\subsection*{What changed}
\begin{itemize}[nosep,leftmargin=1.4em]
   \item \textbf{Unifies} Exp.~2's two-stage search with Exp.~3's normalised
         metric, and \textbf{fixes the folder$\to$task mapping bug}: the two
         earlier scripts had contradicted each other on whether \code{SR} meant
         star\_removal or super\_resolution. The correct pairing
         (\code{SR}$\to$star\_removal, \code{SU}$\to$super\_resolution,
         \code{IR}$\to$image\_restoration) is fixed here, since \emph{every}
         stage-1 baseline depends on pairing each expert with the right folder.
   \item \textbf{Two baselines.} Stage~1 normalises each single-degradation
         folder by its \emph{expert}. Stage~2's combined folders have no single
         expert, so they are normalised by \textbf{$\theta_0$} (the un-merged
         base model) instead.
   \item The grid grows and shifts: $\lambda_{\mathrm{SR}}\in\{0.7,0.8,0.9,1.0\}$,
         $\lambda_{\mathrm{IR}}\in\{0.1,0.3,0.5\}$,
         $\lambda_{\mathrm{SU}}\in\{0.1,0.3,0.5\}$ (36 combos). Crucially
         $\lambda_{\mathrm{SU}}=0$ is \textbf{removed} from the grid, forcing a
         genuine three-task compromise.
\end{itemize}

\subsection*{Rationale}
Combine the best of both prior threads. Normalising the combined folders by
$\theta_0$ answers ``how much does merging improve on doing nothing?''---the
natural reference when no single expert exists for a mixed degradation.

\subsection*{Findings}
The fair metric changes the answer for the first time. Stage~1's winner is
\textbf{\lam{0.9}{0.1}{0.1}} (score $0.7512$) --- \textbf{super-resolution is
finally non-zero}. Stage~2 confirms the same coefficients (score $0.1091$; the
merged model beats $\theta_0$ on the combined folders by large margins, e.g.\
ratios $0.03$--$0.31$). Expert baselines were
$\mathrm{MSE}_{\mathrm{SR}}=0.00094$, $\mathrm{MSE}_{\mathrm{IR}}=0.00244$,
$\mathrm{MSE}_{\mathrm{SU}}=0.00842$, confirming SU's much higher error floor.
Notably the winning merge's \code{SR} ratio is $\sim 1.17$ (slightly worse than
the star-removal expert) while \code{IR} and \code{SU} ratios are \emph{below} 1
--- the normalised metric is now trading tasks off against each other, exactly
as intended.

% ======================================================================
\section{Experiment 5 --- Perceptual Z-Score and Interference Diagnosis}

\subsection*{What changed}
\begin{itemize}[nosep,leftmargin=1.4em]
   \item \textbf{LPIPS} (AlexNet backbone, via the \code{lpips} library) is added
         as a perceptual metric and computed on any folder involving
         super-resolution; \textbf{SSIM is dropped}.
   \item The MSE/baseline ratio is replaced by a \textbf{combined Z-score}: for
         each folder, the relevant metric is standardised
         ($z=(x-\mu)/\sigma$) across all candidate merges and averaged. Folders
         containing \code{SU} are ranked by \textbf{LPIPS}; \code{SR}/\code{IR}
         folders by \textbf{MSE}. Lower is better.
   \item A sampling-free \textbf{deterministic validation loss} is added (the
         v-prediction loss at fixed timesteps with fixed noise), giving a stable
         proxy that does not depend on the stochastic sampler.
   \item Efficiency: DDIM cut to \textbf{10 steps}, evaluation capped at 16
         samples/folder; the full ancestral \code{sample()} method is removed.
         Stage-1 and final models are saved separately.
   \item A dedicated \textbf{diagnostics} cell measures task-vector norms and
         pairwise cosine similarities and compares $\theta_0$, the pure SU expert,
         the winning merge and an ``SU-pushed'' variant at 10 vs.\ 50 DDIM steps.
\end{itemize}

\subsection*{Rationale}
Exp.~4 gave a balanced-looking answer, but visual inspection of super-resolution
outputs did not match the low MSE. \textbf{MSE rewards blur}: a smooth,
detail-free prediction can score well on MSE while looking clearly worse. LPIPS
was introduced precisely to catch this, and the Z-score makes the different
metrics comparable across folders so LPIPS (for SU) and MSE (for SR/IR) can be
combined into one ranking.

\subsection*{Findings}
The Z-score again selects \textbf{\lam{0.9}{0.1}{0.1}}, but the diagnostics
deliver the log's most important result. On the \code{SU} folder:

\begin{center}
\renewcommand{\arraystretch}{1.2}
\footnotesize
\begin{tabular}{@{}lccc@{}}
\toprule
\textbf{Model (on \code{SU})} & \textbf{Det.\ val loss} & \textbf{MSE} & \textbf{LPIPS} \\
\midrule
$\theta_0$ (pure pre-train)          & 0.0835 & 0.00827 & \textbf{0.2303} \\
SU fine-tuned ($\lambda_{\mathrm{SU}}{=}1$) & 0.0880 & 0.00893 & \textbf{0.2203} \\
Winning merge (0.9/0.1/0.1)          & 0.1557 & \textbf{0.00707} & \textbf{0.4080} \\
\bottomrule
\end{tabular}
\end{center}

The winning merge has the \emph{lowest} MSE ($0.00707$) yet by far the
\emph{worst} LPIPS ($0.408$ vs.\ $\sim0.22$ for both the pure base and the pure
SU expert). Two conclusions follow:

\begin{enumerate}[nosep,leftmargin=1.6em]
   \item \textbf{MSE was actively misleading.} The whole MSE-based line of search
         was rewarding a perceptually \emph{worse}, blurrier merged model ---
         vindicating the move to LPIPS.
   \item \textbf{This is task-arithmetic interference, not a sampling artifact.}
         The LPIPS gap persists at 50 DDIM steps ($0.409$), so it is not blur from
         cheap sampling. Since the SU expert \emph{alone} is fine (LPIPS $0.22$)
         but the merge is not, the super-resolution capability is being destroyed
         by interference from the other task vectors when they are summed.
\end{enumerate}

The cosine similarities support this reading: the task vectors are nearly
orthogonal (cos $\approx 0.08$--$0.12$), which is normally favourable for task
arithmetic, but $\lVert\tau_{\mathrm{SU}}\rVert\!\approx\!80.8$ is much larger
than the others, so even a small mixing coefficient on SR/IR perturbs the SU
direction enough to wreck perceptual quality. The notebook's own reading points
the next steps at interference-aware merging --- \textbf{TIES / DARE},
renormalising $\tau_{\mathrm{SU}}$, or using an earlier (fewer-epoch) SU
checkpoint.

% ======================================================================
\section{Overall Trajectory}

The five experiments are best read not as five different models but as a single,
sharpening question: \emph{what does it mean for a merged diffusion model to be
``good'' at three tasks at once?}

\begin{enumerate}[nosep,leftmargin=1.6em]
   \item \textbf{Exp.~1--2} answered it with raw MSE and were repeatedly told that
         the best three-task model is a \emph{two}-task model
         ($\lambda_{\mathrm{SU}}=0$) --- a metric artifact, not a real result.
   \item \textbf{Exp.~3--4} fixed the metric by normalising against per-task
         baselines, and immediately recovered a genuine three-task compromise
         (\lam{0.9}{0.1}{0.1}).
   \item \textbf{Exp.~5} showed that even a good MSE-based answer hides a
         perceptual failure on super-resolution, replaced MSE with LPIPS where it
         matters, and diagnosed the cause as task-vector interference --- turning
         a hyperparameter search into a concrete research direction.
\end{enumerate}

The concrete deliverable stabilised on the coefficient set
\textbf{\lam{0.9}{0.1}{0.1}}, but the durable finding is methodological: for
merging generative restoration models, \textbf{the choice of evaluation metric
matters more than the choice of search algorithm}, and pixel-wise error must be
supplemented by a perceptual metric before any merge can be trusted.

\end{document}

